\documentclass[10pt,journal]{IEEEtran}
% --- Core IEEEtran-friendly packages (keep this list minimal) ---
\usepackage{cite}
\usepackage{amsmath,amssymb}
\usepackage{graphicx}
\usepackage{xcolor}
% Keep IEEEtran's Times setup in a XeTeX preview.
\ifdefined\XeTeXversion
  \usepackage[OT1]{fontenc}
\fi
% Change revisioncolor to black for the clean-text view.
\colorlet{revisioncolor}{blue}
\DeclareRobustCommand{\rev}[1]{\textcolor{revisioncolor}{#1}}
\usepackage[caption=false,font=footnotesize]{subfig}

% Tables
\usepackage{booktabs}
\usepackage{tabularx} % loads array internally
\usepackage{makecell}
\usepackage{multirow}

% Reduce Underfull \hbox warnings in narrow IEEE columns (gentler than \sloppy)
\emergencystretch=1em

% tabularx helper column + thin rule used in Table~\ref{tab:thz_isac_comparison}
\newcolumntype{Y}{>{\centering\arraybackslash}X}
\newcommand{\rowrule}{\specialrule{0.22pt}{0.4pt}{0.4pt}}

% Some TeX installations define these lengths via subfig; define only if missing.
\makeatletter
\@ifundefined{subfigcapskip}{\newlength{\subfigcapskip}}{}
\@ifundefined{subfigbottomskip}{\newlength{\subfigbottomskip}}{}
\makeatother

\title{Self-Interference Reuse for Shared-Aperture Monostatic Full-Duplex THz ISAC}

\author{Sungmin~Cho,~\IEEEmembership{Member,~IEEE,}
        Sang-Rok~Moon,~\IEEEmembership{Member,~IEEE,}
        Sooyeon~Kim,
        Minkyu~Sung,~\IEEEmembership{Member,~IEEE,}
        Seung-Hyun~Cho,~\IEEEmembership{Member,~IEEE,}
        and~Ho-Jin~Song,~\IEEEmembership{Fellow,~IEEE}%
\thanks{This work was supported by the Electronics and Telecommunications
Research Institute (ETRI) grant funded by the Korean government
[26ZH1100, Study on 3D communication technology for hyper-connectivity].
(Corresponding author: Sungmin Cho.)}%
\thanks{S. Cho, S.-R. Moon, S. Kim, M. Sung, and S.-H. Cho are with the Electronics
and Telecommunications Research Institute (ETRI), Daejeon 34129, South
Korea (e-mail: sungmincho@etri.re.kr).}%
\thanks{H.-J. Song is with the Department of Electrical Engineering and
Graduate School of Semiconductor Technology, Pohang University of Science
and Technology (POSTECH), Pohang 37673, South Korea
(e-mail: hojin.song@postech.ac.kr).}%
}

\begin{document}
\maketitle
\begin{abstract}
Terahertz (THz) integrated sensing and communication (ISAC) can
combine high-rate communication with fine ranging. For mobile and vehicular applications, sensing during data transmission helps maintain awareness of changing surroundings by avoiding time-division sensing gaps. Monostatic full-duplex ISAC (FD-ISAC) enables simultaneous downlink
transmission and target-echo reception, but transmitter leakage
causes strong self-interference (SI), particularly in shared-aperture
transceivers. This paper proposes residual-SI reuse, in which passive
\rev{antenna-domain isolation sets the residual leakage level for reuse
as an embedded self-homodyne reference in a square-law receiver.} Because the SI power is independent of target
range, the useful SI--echo cross-beat power follows the conventional
$R^{-4}$ monostatic-radar dependence rather than the $R^{-8}$
echo-only square-law dependence. The SI response from the same
acquisition also enables channel frequency response (CFR)
normalization for ranging. A 280-GHz shared-aperture prototype using
free-running lasers and square-law direct detection \rev{demonstrates concurrent downlink transmission and monostatic
sensing using a fixed wideband passive duplexer and post-detection
CFR baseline removal.} Using a single 20-GHz data-bearing waveform, the system
achieves a 100-Gb/s gross downlink rate over 1~m. \rev{The model identifies a residual-SI power that maximizes sensing
signal-to-interference-plus-noise ratio (SINR) for a fixed echo.
A target displacement of approximately 7.5~mm produces a
corresponding one-range-cell shift in the SI-normalized CFR profile.}
These results support residual-SI reuse as a design approach for
wideband shared-aperture THz FD-ISAC.
\end{abstract}
\begin{IEEEkeywords}
Full-duplex, integrated sensing and communication,
self-interference reuse, square-law detection, terahertz.
\end{IEEEkeywords}

\section{Introduction}

Sixth-generation (6G) networks are expected to combine very-high-rate
connectivity with accurate localization and environmental perception for
applications such as vehicle-to-everything (V2X), autonomous mobility,
and smart manufacturing. Integrated sensing and communication (ISAC)
supports these functions by sharing spectrum, waveforms, hardware, and
signal processing between communication and sensing
\cite{liu_integrated_2022,gonzalez-prelcic_integrated_2024}. The
terahertz (THz) band is particularly attractive for ISAC because its
large available bandwidth can support very-high-rate communication and
fine ranging
\cite{song_terahertz_2022,chaccour_seven_2022,
elbir_Terahertz_2024,zhang_terahertz_2026}. Recent experimental THz ISAC systems have demonstrated these capabilities by focusing on waveform-level integration \cite{jia_unified_2018,lyu_multichannel_2024,shi_310_2026a}.

In mobile and vehicular environments, communication links, target
locations, and beam directions can change rapidly. Time-division ISAC
can therefore leave sensing gaps during communication, delaying the
detection of suddenly appearing or rapidly approaching objects  
\cite{he_nonlinear_2024,zhang_cochannel_2026}. Monostatic full-duplex ISAC (FD-ISAC) addresses these sensing gaps by
enabling simultaneous downlink transmission and monostatic echo
reception at the same node and in the same band
\cite{zhang_inband_2022,duan_frequencydomain_2024,duan_inband_2025,
zhang_cochannel_2026,du_fullduplexbased_2026,du_fullduplex_2026}. Here,
\rev{FD at the monostatic node refers specifically to concurrent downlink
transmission and echo reception. The demonstrated data link is
one-way.} For THz links, where narrow
high-gain beams impose stringent alignment requirements
\cite{sung_practical_2026}, a shared aperture offers the advantages of
a common transmit--receive boresight and reduced antenna duplication
\cite{wu_integrated_2025}.

The main challenge of FD-ISAC is self-interference (SI). In-band
transmitter leakage enters the sensing receiver together with the much
weaker target echo, limiting receiver linearity and dynamic range.
Conventional FD-ISAC therefore suppresses SI in multiple stages,
combining passive antenna-domain isolation with additional RF and digital
SI cancellation (SIC)
\cite{duan_inband_2025,
du_fullduplexbased_2026,du_fullduplex_2026}.
A photonics-assisted in-band FD-ISAC transceiver employed photonic SIC and
achieved 23.6~dB of SI suppression after radar dechirping
\cite{yu_photonicsassisted_2024a}. A recent shared-aperture FD-ISAC
prototype used 30~dB of passive circulator isolation together with RF
and digital SIC, achieving 126~dB of total SI suppression
\cite{du_fullduplex_2026}. \rev{These architectures use dedicated cancellation stages to reduce
leakage before recovering the desired received signals. Cancellation
depth depends on the signal bandwidth, receive functions, and
reference planes used in each demonstration.}

\rev{For a photonic THz downlink with a shared aperture, residual SI
provides a reference field already present at the sensing receiver.
A homodyne receiver with a separately supplied reference would
produce the same delayed beat at equal detector input fields. Here
the duplexer sets the reference through its transmit-to-receive
coupling. This eliminates the separate reference distribution path
in the demonstrated sensing branch. The post-detection SI-suppression
step consists of one complex scalar normalization and a constant
CFR baseline subtraction. It uses the CFR already computed for
ranging and avoids synthesizing a wideband replica of the RF
leakage for this step.}

\rev{At fixed echo power, raising the residual SI initially strengthens
the useful SI--echo beat. SI-induced noise and signal--signal beat
interference (SSBI) eventually limit the gain, giving a
sensing-SINR-maximizing reference power. The prototype implements a
fixed operating point with a passive wideband duplexer. Adjustable
coupling would allow the reference power to be set for a different
operating point.}

The main contributions are summarized as follows.
\begin{itemize}

\item We establish a system model for the proposed shared-aperture
monostatic FD-ISAC architecture and derive the sensing signal-to-interference-plus-noise ratio (SINR) from the
square-law detector output. \rev{We derive the optimum residual-SI power at fixed echo power
and characterize its relation to the receiver input limits and
joint ISAC range. We also formulate SI-normalized CFR ranging
for multiple echoes using an inverse fast Fourier transform
(IFFT).}

\item We experimentally demonstrate the proposed SI-reuse FD-ISAC
architecture at 280~GHz \rev{using a fixed-isolation passive duplexer and digital CFR baseline
removal. The prototype supports} 100-Gb/s gross downlink
communication and monostatic sensing with a 20-GHz bandwidth and a
7.5-mm bandwidth-limited range cell.

\end{itemize}

The remainder of this paper is organized as follows.
Section~\ref{sec:system_model} presents the proposed architecture and
signal model. Section~\ref{sec:experimental_setup} describes the
experimental setup and processing. Section~\ref{sec:experimental_results}
presents the experimental results and comparison.
Section~\ref{sec:conclusion} concludes the paper.



\section{Proposed SI-Reuse Shared-Aperture Monostatic FD-ISAC System}
\label{sec:system_model}

\begin{figure}[!t]
\centering
\includegraphics[width=\columnwidth]{fig/fig1.png}
\caption{Conceptual diagram of the proposed shared-aperture monostatic
full-duplex THz ISAC architecture. LD: laser diode; IM: intensity
modulator; UTC-PD: uni-traveling-carrier photodiode;
LNA: low-noise amplifier; ADC: analog-to-digital converter; DSP: digital
signal processing.}
\label{fig:concept}
\end{figure}

Fig.~\ref{fig:concept} shows the proposed shared-aperture monostatic
FD-ISAC architecture. A duplexing network connects the transmitter and
colocated sensing receiver to a common aperture, providing a common
boresight and passive antenna-domain isolation. \rev{The duplexer reduces the leakage at the receiver input and sets the
residual reference field available for sensing.} At the square-law receiver, its beating with the
target echo provides an embedded self-homodyne sensing reference, while
the same SI response enables same-acquisition CFR normalization.
Section~II-A defines the waveform model, Section~II-B develops the
SI-assisted sensing and CFR normalization, Section~II-C presents the
communication model, and Section~II-D derives the joint ISAC range behavior.

\subsection{Free-Running Photonic Transmitter and Waveform}
\label{sec:tx_signal}

Two free-running lasers separated by $f_{\mathrm c}$ are photomixed in a
uni-traveling-carrier photodiode (UTC-PD), with one optical carrier
modulated by the data waveform. The complex analytic THz signal is
\begin{equation}
x(t)=\big[1+m\,d(t)\big]
e^{j[2\pi(f_{\mathrm c}+\delta f)t+\Delta\theta(t)]},
\label{eq:xt}
\end{equation}
where $d(t)$ is the normalized analytic representation of the
data-bearing intermediate-frequency (IF) waveform, satisfying
$\mathbb{E}[d(t)]=0$ and $\mathbb{E}[|d(t)|^2]=1$, $m$ is the
modulation index, $\delta f$ is the beat-frequency offset, and
$\Delta\theta(t)$ is the relative phase noise. The carrier-to-signal power ratio (CSPR) is
$\mathrm{CSPR}=1/m^2$. With the above normalization, the average
power of the modulated field is $1+m^2$ times its carrier-component
power.

\rev{When the beat-frequency offset $\delta f$ varies slowly within an
acquisition, it contributes a frequency-independent phase to each
delay path. The CFR retains the frequency-dependent phase associated
with that delay.} Orthogonal frequency-division multiplexing (OFDM) enables direct
subcarrier-wise CFR estimation using the known communication symbols.
DFT-spread OFDM (DFT-s-OFDM) is adopted to retain this
frequency-domain processing while reducing the peak-to-average power
ratio (PAPR), allowing more efficient use of the limited THz transmit
power \cite{han_thz_2024}. \rev{In the monostatic architecture, the sensing processor knows the
transmitted data on the occupied subcarriers. The entire data waveform
can therefore serve as the reference for CFR estimation, with a
dedicated sensing pilot optional \cite{keskin_monostatic_2023}.}


\subsection{SI-Assisted Sensing Branch}
\label{sec:sens_branch}

At the sensing low-noise amplifier (LNA) input, SI and the delayed target
echo are superimposed as
\begin{equation}
r(t)=\sqrt{P_{\mathrm{SI}}}\,x(t)
+\sqrt{P_{\mathrm{ec}}(R)}\,x(t-\tau)+n(t),
\label{eq:rt}
\end{equation}
where $\tau=2R/c$ and $n(t)$ is the input-referred receiver noise with
power $N=kT_0F_{\mathrm{LNA}}B$, where $k$ is Boltzmann's constant,
$T_0$ is the reference temperature, $F_{\mathrm{LNA}}$ is the LNA noise
factor, and $B$ is the receiver bandwidth. The residual SI power after
passive antenna-domain isolation is
$P_{\mathrm{SI}}=P_{\mathrm t}10^{-\mathrm{ISO}/10}$, where
$\mathrm{ISO}$ denotes the passive antenna-domain isolation.
\rev{The usable reference power is bounded by the LNA input range and the
square-law range of the detector after LNA gain and interstage loss.}
$P_{\mathrm{SI}}$ is independent of target range, whereas the monostatic
echo power follows \cite{skolnik_introduction_2001}
\begin{equation}
P_{\mathrm{ec}}(R)=
\frac{P_{\mathrm t}G_{\mathrm t}G_{\mathrm r}
\lambda^2\sigma_{\mathrm{eff}}}
{(4\pi)^3R^4L^2},
\label{eq:pec}
\end{equation}
where $G_{\mathrm t}$ and $G_{\mathrm r}$ are the antenna gains,
$\lambda$ is the wavelength, $\sigma_{\mathrm{eff}}$ is the effective
radar cross section (RCS), and $L$ is the one-way front-end loss.
Hereafter, $P_{\mathrm t}$, $P_{\mathrm{SI}}$, and $P_{\mathrm{ec}}$
denote carrier-component powers, with $P_{\mathrm{ec}}(R)$ abbreviated
as $P_{\mathrm{ec}}$ when the range dependence is implicit.

\subsubsection{Square-Law Detection}
\label{sec:sld}

The square-law detector output is proportional to $|r(t)|^2$.
Expanding \eqref{eq:rt} gives
\begin{align}
|r(t)|^2={}&
\underbrace{2\sqrt{P_{\mathrm{SI}}P_{\mathrm{ec}}}
\,\mathrm{Re}\{x(t)x^*(t-\tau)\}}_{\text{SI--echo cross-beat}}
+\underbrace{P_{\mathrm{ec}}|x(t-\tau)|^2}_{\text{echo self-beat}}
\notag\\
&+\underbrace{P_{\mathrm{SI}}|x(t)|^2}_{\text{SI self-beat}}
+\mathcal N(t),
\label{eq:sq_expand}
\end{align}
where
\begin{equation}
\mathcal N(t)
=
2\,\mathrm{Re}\!\left\{
\left[
\sqrt{P_{\mathrm{SI}}}x(t)
+\sqrt{P_{\mathrm{ec}}}x(t-\tau)
\right]n^*(t)
\right\}
+|n(t)|^2.
\label{eq:noise}
\end{equation}

For the waveform in \eqref{eq:xt},
\begin{equation}
|x(t)|^2
=
1+m[d(t)+d^*(t)]+m^2|d(t)|^2.
\label{eq:selfbeat}
\end{equation}
For the self-beat terms, square-law detection removes the common carrier
phase and phase noise. IF bandpass
filtering then extracts the desired beat, which is linear in $m$, while
the second-order term contributes signal--signal beat interference
(SSBI), smaller by a factor of order $m$. A CSPR of 13~dB
($m=0.224$) is therefore used, making $m$ small enough to suppress SSBI
while keeping sufficient sideband power.

The SI--echo cross-beat retains the differential phase
$\phi_\tau(t)=2\pi(f_{\mathrm c}+\delta f)\tau+
\Delta\theta(t)-\Delta\theta(t-\tau)$. The frequency-related term
$2\pi(f_{\mathrm c}+\delta f)\tau$ is constant within each acquisition,
whereas the random variation is determined by the phase-noise difference
over the round-trip delay. For two independent 15-kHz-linewidth lasers,
the resulting 30-kHz beat linewidth gives a root-mean-square (rms)
differential phase variation of $2.9^\circ$ at $R=2$~m.

\rev{After IF filtering, the first-order analytic detector output is}
{\color{revisioncolor}
\begin{equation}
y(t)=2m\big[C_0d(t)+C_\tau d(t-\tau)\big]+n_{\mathrm d}(t).
\label{eq:yt}
\end{equation}}
\rev{Here \(C_0=P_{\mathrm{SI}}+Se^{j\phi_\tau}\) and
\(C_\tau=P_{\mathrm{ec}}+Se^{-j\phi_\tau}\), with
\(S=\sqrt{P_{\mathrm{SI}}P_{\mathrm{ec}}}\). These coefficients
multiply the undelayed and delayed data, respectively.}

\subsubsection{SI-Normalized CFR and Range Estimation}
\label{sec:range_est}

\rev{For \(L_{\mathrm t}\) stationary echoes, the delayed term in
\eqref{eq:yt} becomes a sum over paths. Let \(C_\ell\) and
\(\tau_\ell\) denote the effective detector coefficient and
delay of echo \(\ell\), with the dominant SI at delay zero. The
coefficients are \(C_\ell=a_\ell(\sum_{q=0}^{L_{\mathrm t}}a_q)^*\),
where \(a_0\) is the SI carrier field and \(a_q\) for \(q>0\)
are the echo carrier fields. With negligible inter-block
interference, the occupied-bin model is}
{\color{revisioncolor}
\begin{equation}
\begin{split}
Y_k &\simeq H_kD_k+\mathcal V_k,\\
H_k &=2mA_k\!\left(C_0+
   \sum_{\ell=1}^{L_{\mathrm t}} C_\ell e^{-j2\pi k\Delta f\tau_\ell}
              \right).
\end{split}
\label{eq:Yk}
\end{equation}}
\rev{Here \(D_k\) and \(Y_k\) are the transmit and receive spectra,
\(\Delta f\) is the subcarrier spacing, \(A_k\) is the detector,
IF, and digitizer response, and \(\mathcal V_k\) is receiver noise.
A reference phase at the center frequency is absorbed in \(C_\ell\).
The known \(D_k\) gives a per-bin CFR estimate
\cite{sturm_waveform_2011}. The measurements use the regularized
multi-block estimator in Section~\ref{sec:offline_processing}.}
{\color{revisioncolor}
\begin{equation}
\widehat H_k=Y_k/D_k,\qquad D_k\ne0.
\label{eq:cfr}
\end{equation}}
\rev{With the ideal frequency-resolved SI reference
\(H_{\mathrm{SI},k}=2mA_kC_0\), the expected normalized
response is}
{\color{revisioncolor}
\begin{equation}
\widetilde H_k=\frac{\widehat H_k}{H_{\mathrm{SI},k}}-1
\simeq \sum_{\ell=1}^{L_{\mathrm t}}
\frac{C_\ell}{C_0}e^{-j2\pi k\Delta f\tau_\ell}.
\label{eq:sinorm}
\end{equation}}
\rev{The practical implementation replaces \(H_{\mathrm{SI},k}\) by a
complex scalar estimated across occupied bins. This removes the
common complex scale and the constant baseline, while frequency
selectivity of \(A_k\) can remain.}

\rev{To form the range profile, place the occupied bins \(\mathcal K\)
at their signed FFT indices in an \(N_{\mathrm R}\)-point array
and set the other bins to zero. The IFFT gives}
{\color{revisioncolor}
\begin{equation}
p[n]=\frac{1}{N_{\mathrm R}}\sum_{k\in\mathcal K}
       \widetilde H_k e^{j2\pi kn/N_{\mathrm R}},\qquad
R_n=\frac{cn}{2N_{\mathrm R}\Delta f}.
\label{eq:slope}
\end{equation}}
\rev{After calibration of the SI-path offset, local maxima of
\(\lvert p[n]\rvert^2\) in the search region give echo candidates.
A peak is retained if \(\lvert p[n]\rvert^2>T[n]\) for the chosen
threshold \(T[n]\). Zero padding refines the range grid. The
occupied bandwidth sets the nominal range resolution \(c/(2B)\).
Delayed leakage paths enter \eqref{eq:Yk} as additional terms of
the same delay form and can produce spurious peaks. The scalar
reference is most effective when one leakage path dominates. A
target-off reference can help characterize delayed leakage and
static clutter.}

\subsubsection{Sensing SINR}

\rev{The single-target SINR follows from the delayed term in
\eqref{eq:yt}. For a unit-power analytic IF waveform, the delayed
real-IF power is \(2m^2|C_\tau|^2\). Averaging over carrier phase
with \(\mathbb E[\cos\phi_\tau]=0\) gives
\(2m^2(P_{\mathrm{SI}}P_{\mathrm{ec}}+P_{\mathrm{ec}}^2)\).
The detector-output denominator uses the fixed floor
\(N_{\mathrm{d},0}\), the carrier--noise beat
\(2N(P_{\mathrm{SI}}+P_{\mathrm{ec}})\), and the residual
in-band SSBI. All terms are in input-referred power-product
units. Here \(N\) represents the total input RF noise from the
image bands admitted by the detector and IF filter. Applying
the coherent processing gain gives}
\begin{equation}
\gamma_{\mathrm{sens}}(R)
=
\eta_{\mathrm d}G_{\mathrm p}
\frac{
2m^2\!\left[
P_{\mathrm{SI}}P_{\mathrm{ec}}
+
P_{\mathrm{ec}}^2
\right]
}{
N_{\mathrm{d},0}
+
2N\!\left[
P_{\mathrm{SI}}+P_{\mathrm{ec}}
\right]
+
\kappa m^4P_{\mathrm{SI}}^2
},
\label{eq:sinr_sens}
\end{equation}
where $\eta_{\mathrm d}\leq1$ is the waveform-inversion efficiency and
$\kappa$ is the effective in-band SSBI coefficient after IF filtering
and digital signal processing. Without SI, \eqref{eq:sinr_sens} reduces to
\begin{equation}
\gamma_{\mathrm{sens}}^{\mathrm{w/o\,SI}}(R)
=
\eta_{\mathrm d}G_{\mathrm p}
\frac{
2m^2P_{\mathrm{ec}}^2
}{
N_{\mathrm{d},0}+2NP_{\mathrm{ec}}
},
\label{eq:sinr_nosi}
\end{equation}
where echo-generated SSBI is negligible for a weak echo. Since
$P_{\mathrm{ec}}\propto R^{-4}$, the useful echo-only term scales as
$P_{\mathrm{ec}}^2\propto R^{-8}$.

With SI reuse, square-law detection introduces the SI--echo cross-beat
term $P_{\mathrm{SI}}P_{\mathrm{ec}}$. For a weak echo
($P_{\mathrm{SI}}\gg P_{\mathrm{ec}}$), when detector noise dominates
the denominator, \eqref{eq:sinr_sens} becomes
\begin{equation}
\gamma_{\mathrm{sens}}
\simeq
\eta_{\mathrm d}G_{\mathrm p}
\frac{
2m^2P_{\mathrm{SI}}P_{\mathrm{ec}}
}{
N_{\mathrm{d},0}
}.
\label{eq:sinr_weak}
\end{equation}
Because $P_{\mathrm{SI}}$ is independent of target range, the
SI-assisted sensing SINR scales as $R^{-4}$. Under the same
detector-noise-limited condition, the echo-only SINR scales as
$R^{-8}$. \rev{Residual SI thus supplies the detector with the reference field
needed for the cross-beat.}

\rev{At fixed \(P_{\mathrm{ec}}\), varying \(x=P_{\mathrm{SI}}\) in
\eqref{eq:sinr_sens} gives a derivative with sign
\(N_{\mathrm{d},0}-\kappa m^4(x^2+2P_{\mathrm{ec}}x)\). For
\(\kappa>0\), the unconstrained SINR maximum occurs at}
{\color{revisioncolor}
\begin{equation}
P_{\mathrm{SI}}^\star=
\sqrt{P_{\mathrm{ec}}^2+\frac{N_{\mathrm{d},0}}{\kappa m^4}}
-P_{\mathrm{ec}}.
\label{eq:optimal_si}
\end{equation}}
\rev{The feasible setting is obtained by limiting this value to the
available SI range and the linear input limits. The derivation
holds at fixed echo power. Varying transmit power generally
changes both the echo and SI.}

\subsection{Communication Branch}
\label{sec:comm_branch}

The one-way received communication power is
\begin{equation}
P_{\mathrm{rx}}(R)=
\frac{P_{\mathrm t}G_{\mathrm t}G_{\mathrm r}\lambda^2}
{(4\pi)^2R^2L^2}.
\label{eq:prx}
\end{equation}
\rev{The carrier--sideband beat yields real-IF data power
\(2m^2P_{\mathrm{rx}}^2\). The same power-product model gives the
fixed floor, the carrier--noise beat \(2NP_{\mathrm{rx}}\), and the
in-band SSBI contribution \(\kappa m^4P_{\mathrm{rx}}^2\).
This yields}
\begin{equation}
\gamma_{\mathrm{comm}}(R)=
\frac{2m^2P_{\mathrm{rx}}^2(R)}
{N_{\mathrm{d},0}+2NP_{\mathrm{rx}}(R)
+\kappa m^4P_{\mathrm{rx}}^2(R)}.
\label{eq:sinr_comm}
\end{equation}
Unlike the sensing branch, the one-way communication link contains no
SI--echo cross-beat; its desired component is the carrier--sideband
self-beat, with no coherent sensing-processing gain.


\subsection{ISAC Range Analysis}
\label{sec:isac_range}

The communication and sensing ranges are defined by their respective
SINR thresholds $\gamma_{\mathrm c}$ and $\gamma_{\mathrm{th}}$
\cite{sturm_waveform_2011}. The joint ISAC range is therefore
\begin{equation}
R_{\max}
=
\min\!\left(
R_{\max}^{\mathrm{comm}},
R_{\max}^{\mathrm{sens}}
\right).
\label{eq:rmax}
\end{equation}
Table~\ref{tab:simulation_parameters} summarizes the nominal parameters
used for the closed-form analysis.

\begin{table}[!t]
\caption{Parameters for the Closed-Form Range Analysis}
\label{tab:simulation_parameters}
\centering
\footnotesize
\renewcommand{\arraystretch}{1.07}
\setlength{\tabcolsep}{2.5pt}
\begin{tabular}{@{}
>{\raggedright\arraybackslash}p{0.46\columnwidth}
>{\centering\arraybackslash}p{0.24\columnwidth}
>{\raggedleft\arraybackslash}p{0.24\columnwidth}
@{}}
\toprule
Parameter & Symbol & Value \\
\midrule
Carrier / IF frequency
    & $f_{\mathrm c},f_{\mathrm{IF}}$
    & $280/11$~GHz \\
Waveform / modulation
    & --
    & DFT-s-OFDM / 32-cross-QAM \\
Symbol rate / bandwidth
    & $R_{\mathrm s},B$
    & $20$~GBd / $20$~GHz \\
Gross line rate
    & $R_{\mathrm b}$
    & $100$~Gb/s \\
Bandwidth-limited range cell
    & $\Delta R$
    & $7.5$~mm \\
Transmit power
    & $P_{\mathrm t}$
    & $-10$~dBm \\
CSPR / modulation index
    & $\mathrm{CSPR},m$
    & $13$~dB / $0.224$ \\
Beat linewidth
    & $\Delta\nu_{\mathrm b}$
    & $30$~kHz \\
TX/RX antenna gain
    & $G_{\mathrm t},G_{\mathrm r}$
    & $33$~dBi \\
Front-end loss per pass
    & $L$
    & $2.3$~dB \\
Passive antenna-domain isolation
    & $\mathrm{ISO}$
    & $25$~dB \\
SI power at LNA input
    & $P_{\mathrm{SI}}$
    & $-35$~dBm \\
LNA input compression point
    & $\mathrm{IP}_{\mathrm{1dB}}$
    & $-20$~dBm \\
Effective RCS\textsuperscript{a}
    & $\sigma_{\mathrm{eff}}$
    & $-10$~dBsm \\
LNA gain / noise figure
    & $G_{\mathrm{LNA}},\mathrm{NF}_{\mathrm{LNA}}$
    & $13/8$~dB \\
RF noise power
    & $N$
    & $-63.0$~dBm \\
Detector noise floor\textsuperscript{b}
    & $N_{\mathrm{d},0}$
    & $-89.0$~dB(mW$^2$)\\
Effective in-band SSBI coefficient\textsuperscript{c}
    & $\kappa$
    & $0.06$ \\
Coherent processing gain\textsuperscript{d}
    & $G_{\mathrm p}$
    & $30.1$~dB \\
Waveform-inversion efficiency\textsuperscript{d}
    & $\eta_{\mathrm d}$
    & $1$ \\
Required communication SINR\textsuperscript{e}
    & $\gamma_{\mathrm c}$
    & $15.75$~dB \\
Required sensing SINR\textsuperscript{f}
    & $\gamma_{\mathrm{th}}$
    & $13.2$~dB \\
\bottomrule
\end{tabular}
\vspace{2pt}
\begin{minipage}{\columnwidth}
\scriptsize
\textsuperscript{a}Reference value; varied in
Fig.~\ref{fig:isac_range}.\\
\textsuperscript{b}Equivalent square-law detector noise floor referred
to the LNA input.\\
\textsuperscript{c}Obtained from deterministic 20-GBd DFT-s-OFDM
simulation and rounded to 0.06.\\
\textsuperscript{d}$G_{\mathrm p}$: 1024-point coherent processing;
ideal waveform inversion.\\
\textsuperscript{e}Equivalent threshold in additive white Gaussian noise
(AWGN) for quasi-Gray
32-cross-QAM at a pre-forward-error-correction (FEC) bit error rate
(BER) of $2\times10^{-2}$. \rev{This is the design benchmark for the model, using the stated
labeling.}\\
\textsuperscript{f}Single-bin post-processing threshold for a
nonfluctuating target in AWGN with $P_{\mathrm D}=0.9$ and
$P_{\mathrm{FA}}=10^{-6}$ \cite{skolnik_introduction_2001}.
\end{minipage}
\end{table}

\begin{figure}[!t]
\centering
\includegraphics[width=0.80\columnwidth]{fig/fig2.png}
\caption{Sensing SINR versus LNA-input SI power at $R=1$~m and
$\sigma_{\mathrm{eff}}=-10$~dBsm. The dashed line denotes the
weak-echo approximation in \eqref{eq:sinr_weak}.}
\label{fig:sinr_si}
\end{figure}

\begin{figure}[!t]
\centering
\includegraphics[width=\columnwidth]{fig/fig3.png}\\[-5pt]
\hspace*{0.1\columnwidth}\makebox[0pt]{\footnotesize (a)}%
\hspace*{0.5\columnwidth}\makebox[0pt]{\footnotesize (b)}%
\setlength{\abovecaptionskip}{2pt}
\caption{Communication, sensing, and joint ISAC ranges versus effective
RCS for (a) $P_{\mathrm t}=-10$~dBm and (b) $P_{\mathrm t}=0$~dBm.}
\label{fig:isac_range}
\end{figure}

Using these parameters, Fig.~\ref{fig:sinr_si} shows the sensing SINR
versus residual SI power at $R=1$~m and
$\sigma_{\mathrm{eff}}=-10$~dBsm. As $P_{\mathrm{SI}}$ increases beyond $P_{\mathrm{ec}}$, the
SI--echo cross-beat strengthens and the sensing SINR initially increases
approximately in proportion to $P_{\mathrm{SI}}$, as described by
\eqref{eq:sinr_weak}. At higher SI
levels, SI-induced noise and SSBI limit this gain and eventually
degrade the SINR. The nominal $P_{\mathrm{SI}}=-35$~dBm lies within
the resulting useful SI operating region, \rev{providing approximately 15~dB of modeled sensing-SINR gain over
the echo-only baseline. Equation~\eqref{eq:optimal_si} gives an
unconstrained optimum near \(-25.4\)~dBm for the nominal sweep.
The high-SI region of the curve represents the ideal square-law
model. The achievable operating point is subject to the LNA and
ZBD input limits.}

Fig.~\ref{fig:isac_range} shows how this SINR gain translates into
joint ISAC coverage. At $P_{\mathrm t}=-10$~dBm and
$\sigma_{\mathrm{eff}}=-10$~dBsm, SI reuse increases the sensing range
to 1.6~m, nearly twice that of echo-only square-law sensing. The
SI-assisted sensing range therefore exceeds the 1.0-m communication
range, shifting the joint ISAC coverage from sensing-limited to
communication-limited operation.

This shift also changes the transmit-power scaling of the joint ISAC
range. In the weak-echo, detector-noise-limited regime, the
SI-assisted sensing range and the one-way communication range both
scale as $\sqrt{P_{\mathrm t}}$. \rev{In the model, a 10-dB increase in \(P_{\mathrm t}\) raises the
joint range from about 1.0 to 3.2~m. The corresponding hardware
operating point would require sufficient receiver headroom.}

\section{Experimental Setup}
\label{sec:experimental_setup}

\begin{figure*}[!t]
\centering
\setlength{\subfigcapskip}{-10pt}
\setlength{\subfigbottomskip}{2pt}

\subfloat[]{%
  \includegraphics[width=0.87\textwidth,trim=0 1mm 0 0,clip]
  {fig/fig4a.png}%
  \label{fig:exp_config}%
}\\[-2pt]

\subfloat[]{%
  \includegraphics[width=0.9\textwidth,trim=0 4mm 0 0,clip]
  {fig/fig4b.png}%
  \label{fig:exp_photo}%
}\\[-3pt]

\subfloat[]{%
  \includegraphics[width=0.21\textwidth]{fig/fig4c.png}%
  \label{fig:exp_optical}%
}\hspace{2pt}%
\subfloat[]{%
  \includegraphics[width=0.69\textwidth]{fig/fig4d.png}%
  \label{fig:exp_dsp}%
}

\caption{Experimental implementation of the proposed SI-reuse
shared-aperture monostatic full-duplex THz ISAC system with a
280-GHz carrier: (a) system configuration; (b) testbed photograph;
(c) optical spectrum at the UTC-PD input; and (d) communication
receiver and simultaneous two-channel acquisition, with IF spectra of
the communication (C1) and sensing (C2) channels. MZM: Mach--Zehnder
modulator; AWG: arbitrary waveform generator; OC: optical coupler;
WSS: wavelength-selective switch; EDFA: erbium-doped fiber amplifier;
VOA: variable optical attenuator; ZBD: zero-bias Schottky diode.}
\label{fig:exp}
\end{figure*}

Fig.~\ref{fig:exp}(a) and (b) show the 280-GHz proof-of-concept
testbed for simultaneous downlink communication and monostatic
sensing. The system consists of a free-running photonic THz
transmitter, a shared-aperture monostatic front end, and square-law
receivers based on zero-bias Schottky diodes (ZBDs). A single
data-bearing waveform is used for both functions, while the residual
transmitter leakage through the shared-aperture front end is reused as
the sensing reference.

\subsection{ISAC Waveform and Photonic THz Transmitter}

\rev{A pseudorandom binary sequence of order 15 (PRBS15) was mapped to
32-cross-QAM and organized into DFT-s-OFDM blocks with 1024 occupied
subcarriers. The sensing processor uses the stored transmit
waveform over these same subcarriers.}

An arbitrary waveform generator (AWG, Keysight M8194A) generated the
waveform at symbol rates up to 20~GBd with an 11-GHz IF. At 20~GBd, the 20-GHz occupied bandwidth provides a
100-Gb/s gross line rate and a bandwidth-limited monostatic range cell
of $\Delta R=c/(2B)=7.5$~mm. After IF amplification, the waveform
drove a Mach--Zehnder modulator (MZM, iXBlue MXAN-LN-40) with a
13-dB CSPR.

Two independent free-running laser diodes
(LDs, Pure Photonics PPCL500), each with a 15-kHz Lorentzian
linewidth, were operated at 192.41 and 192.13~THz. LD~\#1 was
modulated by the MZM and combined with the continuous-wave LD~\#2.
A wavelength-selective switch (WSS) selected the two optical carriers
and the lower modulation sideband [Fig.~\ref{fig:exp}(c)], avoiding
double-sideband power fading under laser-frequency drift.

The optical signal was amplified and adjusted by an erbium-doped
fiber amplifier (EDFA) and a variable optical attenuator (VOA) before
photomixing in the UTC-PD (NTT IOD-PMJ-13001). The generated THz
carrier was centered at 280~GHz with a maximum output power of
$-10$~dBm. At 20~GBd, the data-bearing THz spectrum occupies
approximately 259--279~GHz.

\subsection{Shared-Aperture FD-ISAC Front End}
\label{sec:setup_core}

The UTC-PD output was connected to a shared-aperture monostatic front
end comprising a custom dual-circular-polarization (dual-CP) duplexer
and a common corrugated horn. The duplexer converted the transmitted
signal to right-hand circular polarization (RHCP) and routed it to the
horn with a measured insertion loss of 2~dB
\cite{shu_wideband_2019}. Finite isolation between the transmit and
sensing ports produced the residual SI used by the proposed
architecture. \rev{The passive duplexer provides a fixed wideband example of the
reference-setting architecture.}

Two TPX lenses (Thorlabs TPX100) collimated and refocused the 1-m THz
beam. A 25-dBi horn at the remote end served as both the communication
receive antenna and the sensing target. The reflected signal returned
to the common aperture with left-hand circular polarization (LHCP) and
was routed to the sensing port, where it combined with the residual SI
at the LNA input. The experimental link model used a 33-dBi effective
antenna gain and a $-7$-dBsm effective RCS.

The measured passive antenna-domain isolation averaged 23~dB over
259--280~GHz. \rev{For a \(-10\)-dBm transmit signal, this gives an estimated mean
LNA-input SI level of \(-33\)~dBm. With the typical 13-dB LNA gain
and an assumed 1--2~dB interstage loss, the mean ZBD input is
estimated at \(-22\) to \(-21\)~dBm, compared with its reported
\(-20\)-dBm input compression point.}

The superimposed SI and target echo were amplified by a THz LNA
(VDI WR3.4AMP-LN) and square-law detected by a ZBD
(VDI WR3.4ZBD). As described by \eqref{eq:sq_expand} and
\eqref{eq:yt}, the square-law detector generated the SI--echo cross-beat
in the IF band. After DC blocking and IF amplification, the sensing
signal was digitized on channel C2 of a digital storage oscilloscope
(DSO, Keysight UXR0404A) at 256~GSa/s.

The remote communication receiver used the same LNA--ZBD
direct-detection chain and was digitized on channel C1. \rev{The two channels were recorded on the DSO time base
[Fig.~\ref{fig:exp}(d)] with a common transmit waveform reference.}

\begin{figure}[!t]
\centering
\includegraphics[width=1\columnwidth]{fig/fig5.png}
\caption{Measured and signal-to-noise ratio (SNR)-limited error vector
magnitude (EVM) versus
symbol rate over the 1-m link with a 280-GHz carrier, including 32-cross-QAM
constellations at 2 and 20~GBd. \rev{The AWGN threshold is the nominal quasi-Gray benchmark.}}
\label{fig:evm_symbol_rate}
\end{figure}

\subsection{Offline ISAC Processing}
\label{sec:offline_processing}

The stored transmit waveform serves as the common processing reference.\rev{\footnote{The archived records used for reprocessing include an
optional superimposed Zadoff--Chu component with a nominal power weight
of 0.20. It is included in the stored reference.} The ranging waveforms used no cyclic prefix
(\(N_{\mathrm{CP}}=0\)). The DSO data were downconverted and
resampled from 256 to 120~GSa/s. The regularized sensing CFR
uses rectangular FFT windows on the SI-aligned block boundaries.
The 15-GBd record has 15 blocks of 8192 samples, while the checked
20-GBd record has ten blocks of 6144 samples. Each block has 1024
occupied subcarriers. At 1.10~m, the geometric round-trip delay is
about 880 samples, or 10.7\% of the 15-GBd window. Samples at the
start of the echo window can therefore belong to the preceding
transmit block. Equation~\eqref{eq:Yk} neglects this term.}

For communication, the C1 waveform is downconverted, synchronized, and
DFT-s-OFDM demodulated with frequency-domain equalization. \rev{RMS EVM and pre-FEC hard-decision BER are evaluated with the
stored reference and the transmitted bit labeling. The
EVM-derived communication metric is
\(1/\mathrm{EVM}_{\mathrm{rms}}^2\), where
\(\mathrm{EVM}_{\mathrm{rms}}\) is expressed in linear units.}

For sensing, the transmit reference and C2 recording are jointly
processed across the blocks to reduce sensitivity to block-dependent
spectral variations. The regularized CFR is estimated as
\begin{equation}
\widehat H_k=
\frac{\sum_{b=1}^{N_b}Y_{b,k}D_{b,k}^{*}}
{\sum_{b=1}^{N_b}|D_{b,k}|^2+\varepsilon\overline S_{DD}},
\qquad \varepsilon=10^{-3},
\label{eq:cfr_reg}
\end{equation}
where \rev{\(N_b\) is the number of jointly processed
blocks (\(N_b=15\) for the 15-GBd ranging record),}
$D_{b,k}$ and $Y_{b,k}$ are the transmitted and received complex
spectra of block $b$, respectively, and $\overline S_{DD}$ is the
mean accumulated transmit spectral power over the occupied subcarriers.

\rev{The post-detection processor estimates a scalar SI reference
from the complex mean of the occupied CFR bins, divides by it,
and subtracts the unit baseline. This SI-suppression step adds
\(\mathcal O(|\mathcal K|)\) operations per CFR profile. The resulting IFFT profile
from \eqref{eq:slope} is referred to the SI arrival. A one-time
known-range calibration fixes the hardware delay offset.}

\rev{We report the profile peak excess relative to the median
off-target background. In linear power units, \(P_{\mathrm{pk}}\)
is the selected peak and \(P_{\mathrm{fl}}\) is the median floor.}
{\color{revisioncolor}
\begin{equation}
Q_{\mathrm{prof,dB}}=
10\log_{10}\!\left(P_{\mathrm{pk}}/P_{\mathrm{fl}}-1\right),
\quad P_{\mathrm{pk}}>P_{\mathrm{fl}}.
\label{eq:sinr_profile}
\end{equation}}
The target peak is the maximum-power bin within $\pm40$~mm of the
nominal range. The median floor is evaluated over 0.5--2.0~m,
excluding a $\pm2\Delta R$ guard interval around the detected target.
\rev{The off-target bins can include sidelobes and static reflections.
Their median defines a profile statistic. The SINR in
\eqref{eq:sinr_sens} uses the mean modeled interference and
noise power.}

\section{Experimental Results and Discussion}
\label{sec:experimental_results}

\begin{figure}[!t]
\centering
\includegraphics[width=0.72\columnwidth]{fig/fig6.png}
\caption{Matched-filter and SI-normalized CFR range profiles at
15~GBd for a target near 1.10~m. ROI: region of interest.}
\label{fig:cfr_profile}
\end{figure}

\subsection{Communication Performance}

Fig.~\ref{fig:evm_symbol_rate} compares the measured EVM with the
SNR-limited result over the 1-m link. The 2-GBd
constellation shows no noticeable phase-noise-induced rotation, whereas
the 20-GBd constellation exhibits broader noise-like spreading. At high
symbol rates, the small residual degradation from the SNR-limited result
mainly reflects frequency-response nonflatness and other wideband
implementation impairments. Overall, the measured EVM remains close to
the SNR-limited prediction, indicating predominantly SNR-limited
operation.

\rev{At 20~GBd, the 32-cross-QAM DFT-s-OFDM waveform carries a
100-Gb/s gross line rate over approximately 1~m. Reprocessing
the archived record gives an EVM of \(-15.768\)~dB and a
hard-decision BER of 3.875\% (1984 errors in 51\,200 bits)
with its transmitted bit mapping. The \(-15.75\)-dB AWGN
benchmark in Fig.~\ref{fig:evm_symbol_rate} assumes quasi-Gray
labeling.}

\subsection{SI-Normalized CFR Ranging}

Fig.~\ref{fig:cfr_profile} compares matched filtering and SI-normalized
CFR processing for the same 15-GBd acquisition. After one-time range
calibration, both methods estimate the target at approximately 1.098~m,
close to the physical distance of 1.10~m. Both profiles also show a
static response near 0.97~m, attributed to reflection from the TPX lens.

\rev{SI-normalized CFR processing improves the target-to-median-floor
contrast while preserving the estimated target position.} Each subsequent acquisition provides
both the SI reference and target response, enabling range estimation
without a separate reference capture for each measurement.

\begin{figure}[!t]
\centering
\includegraphics[width=\columnwidth]{fig/fig7_profile_and_model.png}
\caption{\rev{15-GBd range analysis. (a) Detector-output model
predictions with a fixed SI level. (b) Profile excess statistics
from one reprocessed capture near 1.10~m, reported on their
own scale.}}
\label{fig:sinr_range}
\end{figure}

\subsection{\rev{SI-Assisted Range Analysis}}

\rev{Fig.~\ref{fig:sinr_range}(a) illustrates the modeled range
dependence at 15~GBd for \(P_{\mathrm t}=-10\)~dBm,
\(P_{\mathrm{SI}}=-33\)~dBm, and
\(\sigma_{\mathrm{eff}}=-7\)~dBsm, with a 29.8-dB effective
sensing processing gain and otherwise nominal model parameters.
The echo-only curve is the model baseline. The fixed-SI cross-beat
sustains a higher predicted sensing SINR over this range.}

\rev{Fig.~\ref{fig:sinr_range}(b) shows the profile statistics
for a 1.10-m record reprocessed with the matched filter and
normalized CFR. The two methods locate the peak at 1.0975~m on
the processing grid. The median-based excess increases from
20.08 to 33.33~dB. Using the mean off-target floor on a common
mask gives 26.10~dB for the CFR peak. In the special case of
exponentially distributed noise-only power, the median is
\(\ln 2\) times the mean. Structured background in a measured
profile can change this ratio. The nominal 13.2-dB single-cell
AWGN threshold applies to the sensing SINR model.}

\begin{figure}[!t]
\centering
\subfloat[]{%
  \includegraphics[width=0.6\columnwidth]{fig/fig8a.png}%
  \label{fig:range_stage}}

\subfloat[]{%
  \includegraphics[width=0.73\columnwidth]{fig/fig8b.png}%
  \label{fig:range_profiles}}
\caption{Resolution-scale displacement measurement:
(a) target translation of approximately 7.5~mm checked with a digital
caliper and
(b) corresponding 20-GBd SI-normalized CFR range profiles showing a
7.5-mm range-peak shift.}
\label{fig:range_resolution}
\end{figure}

\begin{table*}[!t]
\caption{Representative millimeter-wave, terahertz,
and full-duplex ISAC demonstrations}
\label{tab:thz_isac_comparison}
\centering
\footnotesize
\setlength{\tabcolsep}{1.0pt}
\renewcommand{\arraystretch}{1.08}

\renewcommand{\tabularxcolumn}[1]{m{#1}}

\begin{tabularx}{\textwidth}{@{}
>{\raggedright\arraybackslash}m{0.105\textwidth}
*{7}{Y}
@{}}
\toprule
&
\makecell{JSTQE'23\\\cite{bai_photonicsassisted_2023}}
&
\makecell{JLT'24\\\cite{lyu_multichannel_2024}}
&
\makecell{JLT'26\\\cite{shi_310_2026a}}
&
\makecell{OE'24\\\cite{yu_photonicsassisted_2024a}}
&
\makecell{IoTJ'25\\\cite{duan_inband_2025}}
&
\makecell{WCM'26\\\cite{du_fullduplex_2026}}
&
\textbf{This work}
\\
\midrule

Sensing mode
& Monostatic
& Bistatic
& Bistatic
& Monostatic
& Monostatic
& Monostatic
& \textbf{Monostatic}
\\
\rowrule

TX/RX aperture
& Separate
& Separate
& Separate
& Separate
& Separate
& Shared
& \textbf{Shared}
\\
\rowrule

SI strategy
& --
& --
& --
& \makecell{Photonic SIC\\(23.6~dB)}
& \makecell{Passive + spatial\\+ digital SIC}
& \makecell{Passive + RF\\+ digital SIC\\(126~dB)}
& \makecell{\textbf{Passive isolation}\\\textbf{+ SI reuse}\\
\rev{+ CFR baseline}\\\rev{removal}}
\\
\rowrule

Sensing RX
& \makecell{Photonic dechirp\\+ PD}
& \makecell{Schottky mixer\\+ LO}
& \makecell{FMD direct\\detection}
& \makecell{PD + mixer after\\photonic SIC}
& \makecell{$16{\times}16$\\phased-array}
& \makecell{USRP B210\\SDR RX}
& \makecell{\textbf{ZBD direct}\\\textbf{detection}}
\\
\rowrule

Ranging
& \makecell{Dechirp\\+ fusion}
& \makecell{Matched filter\\+ fusion}
& \makecell{Weighted log-SNR\\+ path-loss model}
& Dechirp
& \makecell{OFDM range--\\Doppler}
& \makecell{Dechirp\\+ DFT}
& \makecell{\textbf{SI-normalized}\\\textbf{CFR}}
\\

\midrule\midrule

Carrier
& 38/52~GHz
& \makecell{250/275/\\300~GHz}
& 300~GHz
& 10~GHz
& 26.8~GHz
& 1.74~GHz
& \textbf{280~GHz}
\\
\rowrule

Waveform
& CE-LFM-OFDM
& LFM-16QAM
& PCS-OFDM
& CE-LFM-OFDM
& 16-QAM OFDM
& OFDM--chirp
& \makecell{\textbf{32-cross-QAM}\\\textbf{DFT-s-OFDM}}
\\
\rowrule

Bandwidth
& $4~\mathrm{GHz}\times2$
& $10~\mathrm{GHz}\times3$
& 33~GHz
& \makecell{2~GHz total\\(1.5~GHz sensing)}
& 491.52~MHz
& 40~MHz
& \textbf{20~GHz}
\\
\rowrule

Gross data rate
& 16~Gb/s
& 120~Gb/s
& 310.4~Gb/s\textsuperscript{a}
& 3.0~Gb/s\textsuperscript{b}
& N/R
& N/R
& \textbf{100~Gb/s}
\\
\rowrule

Range resolution
& 8.6~mm\textsuperscript{c}
& 2.5~mm\textsuperscript{c}
& N/R
& 100~mm
& 305.2~mm
& N/R
& \textbf{7.5~mm}\textsuperscript{d}
\\
\rowrule

Comm. distance
& 1.2~m
& \multirow{2}{*}{0.5~m}
& 0.3~m
& \multirow{2}{*}{1.3~m}
& N/R
& 40~m
& \multirow{2}{*}{\textbf{1.0~m}}
\\

\cmidrule(lr){1-2}
\cmidrule(lr){4-4}
\cmidrule(lr){6-7}

Sensing distance
& 3.2~m
&
& 0.3--1.0~m
&
& $\sim$85--145~m
& 5--20~m
&
\\

\bottomrule
\end{tabularx}

\vspace{1.5pt}
\parbox{\textwidth}{\scriptsize
\textsuperscript{a}Gross line rate obtained using
dual-polarization 33-GBd PCS-64-QAM and satisfying the NGMI
criterion at 0.3~m.
\textsuperscript{b}Calculated from the reported 64-QAM and
500-MBd CE-OFDM.
\textsuperscript{c}Obtained after frequency-channel fusion.
\textsuperscript{d}Bandwidth-limited monostatic range cell,
assessed through single-target displacement using range-profile
analysis inspired by \cite{lyu_multichannel_2024}.
``--'' denotes not applicable. N/R: not reported.
CE: constant envelope; DFT: discrete Fourier transform; FMD:
Fermi-level-managed barrier diode; LFM: linear frequency modulation;
NGMI: normalized generalized mutual information; PCS: probabilistic
constellation shaping; SDR: software-defined radio; SIC:
self-interference cancellation; USRP: universal software radio peripheral.
}
\end{table*}

\subsection{Resolution-Scale Displacement Measurement}

In the present shared-aperture FD-ISAC setup, the communication
receiver also serves as the sensing target. Displacement sensing is
therefore evaluated by translating this receiver along the range
direction and tracking the SI-normalized CFR peak. The symbol rate
was increased to 20~GBd, providing a 20-GHz occupied bandwidth and
a bandwidth-limited monostatic range cell of 7.5~mm.

The applied displacement of approximately 7.5~mm was independently verified
using a digital caliper [Fig.~\ref{fig:range_resolution}(a)].
The corresponding CFR peak shifts from 1012.5 to 1005.0~mm
[Fig.~\ref{fig:range_resolution}(b)], yielding a 7.5-mm shift
that closely matches the physical displacement. \rev{This record demonstrates displacement on the scale of one
bandwidth-limited range cell using the data-bearing
communication waveform and the acquired SI reference.}

\subsection{Comparison and Future Research Directions}
\label{sec:discussion}

Table~\ref{tab:thz_isac_comparison} compares representative
millimeter-wave/THz and FD-ISAC demonstrations. Previous
waveform-integrated ISAC systems have achieved communication rates up
to 310.4~Gb/s and millimeter-scale ranging
\cite{lyu_multichannel_2024,shi_310_2026a}, but do not explicitly
address SI-limited same-band monostatic FD operation. Recent FD-ISAC
prototypes instead suppress SI using photonic SIC
\cite{yu_photonicsassisted_2024a}, passive and spatial isolation
followed by digital SIC \cite{duan_inband_2025}, or cascaded passive,
RF, and digital SIC \cite{du_fullduplex_2026}. In particular, the
shared-aperture prototype in \cite{du_fullduplex_2026} combines
30~dB of passive circulator isolation with additional RF and digital
SIC to achieve 126~dB of total cancellation. However, the compared
FD-ISAC prototypes operate below 30~GHz with sensing bandwidths up to
1.5~GHz, leaving wideband THz shared-aperture FD-ISAC largely
unexplored.

\rev{The 280-GHz prototype uses a fixed passive duplexer to place
residual SI within a useful reference regime. The local CFR
baseline removal acts after the SI--echo beat has been formed.
Its processing scope differs from the cascaded SI suppression
in \cite{du_fullduplex_2026}, which reports 126~dB for its
configuration. Table~\ref{tab:thz_isac_comparison} places the
two architectures in their respective bandwidth and receive
contexts.}

The primary focus of this work is SI-reuse operation in a
shared-aperture FD-ISAC architecture. To evaluate this principle, we performed a simplified proof-of-concept
experiment under controlled indoor conditions, focusing on static,
one-dimensional ranging near 1~m. The demonstrated range is
primarily limited by the current THz
link budget, including available THz transmit power and interconnect
loss. \rev{Higher transmit power, antenna gain, and detector responsivity
can extend the modeled range subject to receiver headroom.} Chip-level
integration and optimized packaging can further reduce interconnect
loss and enable a smaller form factor.

Future research directions include dynamic and multidimensional
sensing. The SI--echo cross-beat preserves target-dependent Doppler information of a
moving target, making range--Doppler sensing a natural extension of
the present static experiment. Such processing will require
compensation for inter-acquisition frequency and phase drift from the
free-running transmitter. Multi-antenna or multiple-input multiple-output
(MIMO) implementations could
further extend the architecture to angular estimation and
multidimensional localization. The colocated monostatic architecture
can facilitate transmit--receive timing synchronization, while direct
detection avoids separate coherent THz local-oscillator distribution
and locking. Multichannel operation would nevertheless require
channel-to-channel calibration and management of SI coupling across
the array. Finally, the present work realizes FD-ISAC in the
transmit--echo sense rather than bidirectional FD communication.
Extending the architecture to simultaneous uplink and downlink
communication with sensing would additionally require separation of
the uplink signal, residual SI, and target echoes within the available
receiver dynamic range.

\section{Conclusion}
\label{sec:conclusion}

This paper proposed and experimentally demonstrated residual-SI reuse
for shared-aperture monostatic THz FD-ISAC. Passive antenna-domain
\rev{isolation sets the residual leakage level available as an
embedded self-homodyne sensing reference.} The
square-law model shows that useful SI--echo cross-beat power follows
the conventional $R^{-4}$ monostatic-radar dependence, compared with
the $R^{-8}$ echo-only dependence. The SI response also enables
same-acquisition CFR normalization for ranging. The analysis further
\rev{gives the SINR-maximizing residual reference power at fixed echo
power, with the feasible setting constrained by the receiver
input limits.}

A 280-GHz prototype using free-running lasers and square-law direct
detection demonstrated simultaneous downlink communication and
monostatic sensing \rev{using a passive duplexer and post-detection CFR baseline removal.} Using a single 20-GHz 32-cross-QAM DFT-s-OFDM
waveform, the system achieved a 100-Gb/s gross downlink rate over
1~m, while a 7.5-mm target displacement produced the corresponding
SI-normalized CFR peak shift. Together, the analysis and measurements
\rev{support residual-SI reuse for concurrent downlink transmission
and monostatic ranging in a wideband shared-aperture photonic
receiver.}


\bibliographystyle{IEEEtran}
\bibliography{ref_fin}
\end{document}
